% Frugal AI seminar, live demo: the same small model with and without retrieval (FrugalAI/code/demo). Screenshots are the fallback for the live run.
% CheatSheet-safe: no floats. Five fixed questions, one picture each; picked after a trial run, so the setup frame says they are examples.

\section[Demo]{Live Demo: No Internet}

%%%%%%%%%%%%%%%%%%%%%%%%%%%%%%%%%%%%%%%%%%%%%%%%%%%
\begin{frame}[fragile]\frametitle{}
\begin{center}
{\Large Live Demo: No Internet}

\medskip
{\small (One small model on one laptop, Wi-Fi switched off)}
\end{center}
\end{frame}

%%%%%%%%%%%%%%%%%%%%%%%%%%%%%%%%%%%%%%%%%%%%%%%%%%%
\begin{frame}[fragile]\frametitle{The Setup}
\begin{itemize}
\item One small model, Qwen3 1.7B, runs on this laptop through LM Studio.
\item Each question goes to it twice: alone, and with RAG, which first fetches the matching parts of a 39-page syllabus PDF.
\item The five questions were picked as clear examples after a trial run, not as a random sample. Over 110 questions, the model with RAG answered 67\% correctly (one run, one machine).
\item The questions were drafted with Claude's help and the author has not reviewed them yet.
\end{itemize}
{\tiny (Ref: FrugalAI code/demo)}
\end{frame}

%%%%%%%%%%%%%%%%%%%%%%%%%%%%%%%%%%%%%%%%%%%%%%%%%%%
\begin{frame}[fragile]\frametitle{A Number the Model Cannot Guess}
\begin{center}
\adjustbox{max width=0.97\linewidth, max totalheight=0.80\textheight}{\includegraphics{frugal_rag_q1}}
\end{center}
\end{frame}

%%%%%%%%%%%%%%%%%%%%%%%%%%%%%%%%%%%%%%%%%%%%%%%%%%%
\begin{frame}[fragile]\frametitle{An Invented Answer Against an Honest One}
\begin{center}
\adjustbox{max width=0.97\linewidth, max totalheight=0.80\textheight}{\includegraphics{frugal_rag_q2}}
\end{center}
\end{frame}

%%%%%%%%%%%%%%%%%%%%%%%%%%%%%%%%%%%%%%%%%%%%%%%%%%%
\begin{frame}[fragile]\frametitle{A Title the Model Made Up}
\begin{center}
\adjustbox{max width=0.97\linewidth, max totalheight=0.80\textheight}{\includegraphics{frugal_rag_q3}}
\end{center}
\end{frame}

%%%%%%%%%%%%%%%%%%%%%%%%%%%%%%%%%%%%%%%%%%%%%%%%%%%
\begin{frame}[fragile]\frametitle{Finding a Course by Its Name}
\begin{center}
\adjustbox{max width=0.97\linewidth, max totalheight=0.80\textheight}{\includegraphics{frugal_rag_q4}}
\end{center}
\end{frame}

%%%%%%%%%%%%%%%%%%%%%%%%%%%%%%%%%%%%%%%%%%%%%%%%%%%
\begin{frame}[fragile]\frametitle{Credits of a Course, Not a TV Show}
\begin{center}
\adjustbox{max width=0.97\linewidth, max totalheight=0.80\textheight}{\includegraphics{frugal_rag_q5}}
\end{center}
\end{frame}

%%%%%%%%%%%%%%%%%%%%%%%%%%%%%%%%%%%%%%%%%%%%%%%%%%%
\begin{frame}[fragile]\frametitle{A Second Case: Is It Already Covered?}
\begin{itemize}
\item A concept-intake workflow on the author's own quantum FAQ notes: is this concept already covered, and where else does it appear?
\item Plain code does the counting, the searching and the deciding. The small model writes only a draft, and only for a new concept.
\item Measured on 32 test questions from the notes: the right entry came first for all 20 reworded ones, and covered against new was right 91\% of the time (one run).
\item The pictures show real entry names from the author's notes. The five questions were picked as examples.
\end{itemize}
{\tiny (Ref: FrugalAI code/workflow-agent; the numbers come from private notes and cannot be reproduced from the repository alone)}
\end{frame}

%%%%%%%%%%%%%%%%%%%%%%%%%%%%%%%%%%%%%%%%%%%%%%%%%%%
\begin{frame}[fragile]\frametitle{The Workflow, Drawn from the Code}
\begin{center}
\adjustbox{max width=0.97\linewidth, max totalheight=0.70\textheight}{\includegraphics{frugal_faq_graph}}

{\tiny Grey boxes are plain code; the orange box is the small model. Drawn from the compiled LangGraph.}
\end{center}
\end{frame}

%%%%%%%%%%%%%%%%%%%%%%%%%%%%%%%%%%%%%%%%%%%%%%%%%%%
\begin{frame}[fragile]\frametitle{The Path of a New Concept}
\begin{center}
\adjustbox{max width=0.97\linewidth, max totalheight=0.69\textheight}{\includegraphics{frugal_faq_path}}

{\scriptsize One real run: three drafts, each failing the house-rules check (answer too short). Nothing is written to the notes.}
\end{center}
\end{frame}

%%%%%%%%%%%%%%%%%%%%%%%%%%%%%%%%%%%%%%%%%%%%%%%%%%%
\begin{frame}[fragile]\frametitle{Already in the FAQ: Parallelism}
\begin{center}
\adjustbox{max width=0.97\linewidth, max totalheight=0.80\textheight}{\includegraphics{frugal_faq_q1}}
\end{center}
\end{frame}

%%%%%%%%%%%%%%%%%%%%%%%%%%%%%%%%%%%%%%%%%%%%%%%%%%%
\begin{frame}[fragile]\frametitle{Already in the FAQ: Faster-Than-Light Signals}
\begin{center}
\adjustbox{max width=0.97\linewidth, max totalheight=0.80\textheight}{\includegraphics{frugal_faq_q2}}
\end{center}
\end{frame}

%%%%%%%%%%%%%%%%%%%%%%%%%%%%%%%%%%%%%%%%%%%%%%%%%%%
\begin{frame}[fragile]\frametitle{Already in the FAQ: Encryption}
\begin{center}
\adjustbox{max width=0.97\linewidth, max totalheight=0.80\textheight}{\includegraphics{frugal_faq_q3}}
\end{center}
\end{frame}

%%%%%%%%%%%%%%%%%%%%%%%%%%%%%%%%%%%%%%%%%%%%%%%%%%%
\begin{frame}[fragile]\frametitle{A New Concept: Surface Code}
\begin{center}
\adjustbox{max width=0.97\linewidth, max totalheight=0.80\textheight}{\includegraphics{frugal_faq_q4}}
\end{center}
\end{frame}

%%%%%%%%%%%%%%%%%%%%%%%%%%%%%%%%%%%%%%%%%%%%%%%%%%%
\begin{frame}[fragile]\frametitle{A New Concept: No-Cloning}
\begin{center}
\adjustbox{max width=0.97\linewidth, max totalheight=0.80\textheight}{\includegraphics{frugal_faq_q5}}
\end{center}
\end{frame}
